\documentclass[10pt,a4paper]{amsart}
\usepackage[margin=19mm]{geometry}
\usepackage{amsmath,amssymb,mathtools}
\usepackage[scr=rsfs,cal=euler]{mathalfa}
\usepackage{xfrac}
\usepackage[unicode,hidelinks,bookmarks=false]{hyperref}
\newcommand{\ie}{i.e.\ }
\newcommand{\cf}{cf.\ }
\newcommand{\resp}{respectively\ }
\newcommand{\Subsection}{\subsection}
\providecommand{\nobreakdash}{\nobreak\hbox{-}}
\setlength{\emergencystretch}{2em}
\multlinegap0pt
\usepackage{amssymb}
\usepackage[shortlabels]{enumitem}
\usepackage{bm}
\usepackage{algorithm}\usepackage{algpseudocode}
\usepackage{tikz}
\newtheorem{lemma}{Lemma}
\newtheorem{theorem}{Theorem}
\title{Exact Formulas for Coprime Representations of Even Integers Avoiding a Prime}
\author{Andres M. Salazar}
\date{Source: arXiv:2604.02386v1 (CC BY 4.0)}
\begin{document}
\maketitle

\noindent\emph{Source and licence.} Exact Formulas for Coprime Representations of Even Integers Avoiding a Prime. Version arXiv:2604.02386v1 (CC BY 4.0), distributed by arXiv under the Creative Commons Attribution 4.0 International licence, http://creativecommons.org/licenses/by/4.0/, as recorded in the arXiv licence metadata for this exact version and shown on the abstract page https://arxiv.org/abs/2604.02386v1. Full licence notice: \texttt{licenses/arxiv-2604.02386-CC-BY-4.0.txt}.\par\smallskip

\noindent\textit{Editorial scope.} Counted unit: the article's theorem thm:min\_sol\_afrc (source0.tex lines 536--568). The article's complete original proof of it is delivered with it (source0.tex lines 570--749, one \texttt{proof} environment of 180 lines). No mathematics is written, repaired or re-derived: the statement and the proof are the article's own lines, in the article's own notation, and no retained line is substituted or re-flowed.  Retained uncounted as the source-local context the proof needs: lines 348--350; lines 432--436; lines 456--482.  Nothing was omitted from any retained span, so the package carries no omission record.  No retained line contains a citation, so the wrapper carries no bibliography and no bibliography entry.  No retained line contains a figure environment, an image-inclusion command or drawing code, so no image is delivered and the package declares no asset dependency.  Editorial formatting only, in addition: an AMS \texttt{amsart} preamble that restates the article's own declarations and macros used by the retained lines, namely the environments \texttt{lemma}, \texttt{theorem} on separate counters, together with the standard packages those lines need.  The retained environments print in this document as Lemma 1, Theorem 1 in order of appearance, whereas the article prints the same items with the numbering of its own full document.  The complete native source of the article as distributed in the arXiv e-print for this exact version is bundled unchanged as \texttt{sources/197733/source0.tex}.
\begin{equation}\label{eq:general}
\delta_k(ax+b) = c,
\end{equation}

\begin{lemma}[Uniqueness under coprimality]\label{lema:unicidad_ep}
Let $k \ge 2$, $a,b \in \mathbb{Z}$, $c \in \{0,\dots,k-1\}$, and suppose
$\gcd(a,k)=1$. Then equation~\eqref{eq:general} has exactly one solution
in $\{0,\dots,k-1\}$, which coincides with the minimal solution.
\end{lemma}

\begin{lemma}\label{lem_eu}
Let $k \ge 2$ and $a_0 \in \mathbb{Z}$ with $2 \le a_0 < k$. Define
\[
a^{(0)} = k,
\quad
a^{(1)} = a_0,
\quad
a^{(n)} =
\delta_{a^{(n-1)}}\!\bigl(a^{(n-2)}\bigr),
\qquad n \ge 2.
\]
Then:
\begin{enumerate}
\item The sequence satisfies
\[
0 \le a^{(n)} < a^{(n-1)}
\quad \text{whenever } a^{(n-1)} \ne 0.
\]
\item There exists a unique index $r \ge 1$ such that
\[
a^{(r)} = \gcd(a_0,k),
\quad
a^{(r+1)} = 0.
\]
\item The algorithm terminates in $O(\log\min\{k,a_0\})$ steps.
\end{enumerate}
\end{lemma}

\begin{theorem}\label{thm:min_sol_afrc}
Let
\[
k \ge 2,
\quad
\gcd(a_0,k)=1,
\quad
b_0 \in \{0,\dots,k-1\}.
\]
Define sequences
\[
a^{(0)}=k,\quad a^{(1)}=a_0,\quad
a^{(n)}=\delta_{a^{(n-1)}}(a^{(n-2)}),
\]
and
\[
b^{(1)}=b_0,\quad
b^{(n)}=a^{(n-1)}-\delta_{a^{(n-1)}}(b^{(n-1)}).
\]
Let $r$ be the unique index with $a^{(r)}=1$ (cf.\ Lemma~\ref{lem_eu}).
Then the unique solution in $\{0,\dots,k-1\}$ of
\[
\delta_k(a_0 x)=b_0
\]
is
\[
x =
\delta_k\!\left(
k\sum_{j=0}^{r-1}
\frac{b^{(j+1)}}{a^{(j)}a^{(j+1)}}
\right).
\]
\end{theorem}

\begin{proof}
If $b_0 = 0$, then $x = 0$ is the unique solution in $\{0,\dots,k-1\}$,
and the formula gives $\delta_k(0) = 0$, consistently.
Henceforth assume $b_0 \ge 1$.

Since $\gcd(a_0,k)=1$, by Lemma~\ref{lema:unicidad_ep} the equation
admits a unique solution in $\{0,\dots,k-1\}$.

\medskip
\textit{Step 1: Construction of the auxiliary system.}

Set $x_1 := x$ and write the equation as
\[
\delta_{a^{(0)}}(a^{(1)} x_1) = b^{(1)}.
\]
By definition of $\delta_{a^{(0)}}$, there exists an integer $x_2$ such
that
\[
a^{(1)} x_1 = a^{(0)} x_2 + b^{(1)}.
\tag{$R_2$}
\]

We claim that for each $n = 2, \dots, r$ there exists an integer $x_n$
satisfying
\[
a^{(n-1)} x_{n-1} = a^{(n-2)} x_n + b^{(n-1)}.
\tag{$R_n$}
\]

We proceed by induction on $n$.

\textit{Base case} $n = 2$: this is $(R_2)$ already established.

\textit{Inductive step}: assume $(R_n)$ holds for some $2 \le n < r$,
so
\[
a^{(n-1)} x_{n-1} = a^{(n-2)} x_n + b^{(n-1)}.
\]
Since $a^{(n-2)} x_n + b^{(n-1)}$ is divisible by $a^{(n-1)}$,
applying $\delta_{a^{(n-1)}}$ to both sides and using additive
compatibility gives
\[
0 = \delta_{a^{(n-1)}}\!\left(a^{(n-2)} x_n + b^{(n-1)}\right)
  = \delta_{a^{(n-1)}}\!\left(\delta_{a^{(n-1)}}(a^{(n-2)})\, x_n
    + \delta_{a^{(n-1)}}(b^{(n-1)})\right).
\]
Setting $a^{(n)} := \delta_{a^{(n-1)}}(a^{(n-2)})$, this becomes
\[
\delta_{a^{(n-1)}}(a^{(n)} x_n) = a^{(n-1)} - \delta_{a^{(n-1)}}(b^{(n-1)})
= b^{(n)}.
\]
By definition of $\delta_{a^{(n-1)}}$, there exists an integer $x_{n+1}$
such that
\[
a^{(n)} x_n = a^{(n-1)} x_{n+1} + b^{(n)},
\]
which is $(R_{n+1})$. This completes the induction.

\medskip
\textit{Step 2: Base case of back-substitution.}

At step $n = r$, relation $(R_r)$ reads
\[
a^{(r-1)} x_{r-1} = a^{(r-2)} x_r + b^{(r-1)}.
\]
Since $a^{(r)} = \delta_{a^{(r-1)}}(a^{(r-2)}) = 1$, the equation
$\delta_{a^{(r-1)}}(a^{(r)} x_{r-1}) = b^{(r)}$ reduces to
$\delta_{a^{(r-1)}}(x_{r-1}) = b^{(r)}$,
so we may set
\[
x_r := b^{(r)} = a^{(r-1)} - \delta_{a^{(r-1)}}(b^{(r-1)})
\in \{0, \dots, a^{(r-1)}-1\} \subset \mathbb{Z}_{\ge 0}.
\]

\medskip
\textit{Step 3: Back-substitution, integrality, and the closed formula.}

We prove by downward induction on $\ell = r, r-1, \dots, 1$ that
\[
x_\ell
=
\sum_{j=\ell-1}^{r-1}
b^{(j+1)}
\prod_{i=\ell-1}^{j-1}
\frac{a^{(i)}}{a^{(i+1)}}.
\tag{$F_\ell$}
\]

\textit{Base case} $\ell = r$: the sum reduces to the single term
$j = r-1$, giving
\[
x_r = b^{(r)} \cdot \prod_{i=r-1}^{r-2}(\cdots) = b^{(r)} \cdot 1 = b^{(r)},
\]
where the empty product (lower index exceeds upper index) equals $1$.
This agrees with Step~2.

\textit{Inductive step}: assume $(F_\ell)$ holds for some
$2 \le \ell \le r$.
Rearranging $(R_\ell)$,
\[
x_{\ell-1}
=
\frac{a^{(\ell-2)} x_\ell + b^{(\ell-1)}}{a^{(\ell-1)}}.
\]
Substituting $(F_\ell)$,
\begin{align*}
x_{\ell-1}
&=
\frac{a^{(\ell-2)}}{a^{(\ell-1)}}
\sum_{j=\ell-1}^{r-1}
b^{(j+1)}
\prod_{i=\ell-1}^{j-1}
\frac{a^{(i)}}{a^{(i+1)}}
+
\frac{b^{(\ell-1)}}{a^{(\ell-1)}}
\\
&=
\sum_{j=\ell-1}^{r-1}
b^{(j+1)}
\prod_{i=\ell-2}^{j-1}
\frac{a^{(i)}}{a^{(i+1)}}
+
b^{(\ell-1)}
\prod_{i=\ell-2}^{\ell-3}
\frac{a^{(i)}}{a^{(i+1)}}
\\
&=
\sum_{j=\ell-2}^{r-1}
b^{(j+1)}
\prod_{i=\ell-2}^{j-1}
\frac{a^{(i)}}{a^{(i+1)}},
\end{align*}
which is $(F_{\ell-1})$. This completes the induction.

Setting $\ell = 1$ in $(F_1)$ and noting that $a^{(0)} = k$, we obtain
\[
x_1
=
\sum_{j=0}^{r-1}
b^{(j+1)}
\prod_{i=0}^{j-1}
\frac{a^{(i)}}{a^{(i+1)}}.
\]
The product telescopes:
\[
\prod_{i=0}^{j-1}
\frac{a^{(i)}}{a^{(i+1)}}
=
\frac{a^{(0)}}{a^{(j)}}
=
\frac{k}{a^{(j)}},
\qquad j \ge 1,
\]
and for $j = 0$ the empty product equals $1 = k/a^{(0)} = k/k$,
so the formula $k/a^{(j)}$ holds uniformly for all $0 \le j \le r-1$.
Therefore
\[
x_1
=
k \sum_{j=0}^{r-1}
\frac{b^{(j+1)}}{a^{(j)} a^{(j+1)}}
\;\in\; \mathbb{Z},
\]
and by construction $\delta_k(a_0 x_1) = b_0$.

\medskip
\textit{Step 4: Minimality.}

By Lemma~\ref{lema:unicidad_ep} there is exactly one solution in
$\{0,\dots,k-1\}$, so
\[
x = \delta_k(x_1)
=
\delta_k\!\left(
k\sum_{j=0}^{r-1}
\frac{b^{(j+1)}}{a^{(j)}a^{(j+1)}}
\right)
\]
is that unique solution, which coincides with the minimal solution.
\end{proof}

\end{document}
